% page 513 of the facsimile (image p-512 of the scan)
% block 160.0 x 242.0 mm ; left margin 42.0 mm
\pgc{20.320mm}{0.000mm}{
\l{\cel{50}505.}
\l{}
\l{\sou{rufa}. Reda kelke-flava. - L.}
\l{}
\l{\sou{rugo}. Plisuro di la pelo sur la fronto, la vizajo, la manui,}
\l{\cel{3}quan produktas, che la oldi, la mult-evozeso. - (\sou{metaf}.:}
\l{\cel{3}Sulketo sur surfaco di irgo). - FIS.}
\l{}
\l{\sou{ruinar}. (\sou{trans.}) I. Igar (konstrukturo) krular. - II. Igar}
\l{\cel{3}(persono) ne plus havar definitive sua havaji importoza.}
\l{\cel{3}- DEFIRS.}
\l{}
\l{\sou{ruko}.\cel{2}Stipo di kano, bestono di qua la extremajo supra garnis-}
\l{\cel{3}esis ye lino, kanabo, kotono, e c. por filifo. - DEIS.}
\l{}
\l{\sou{ruktar}. (\sou{netrans.}). (\sou{aludante la sto}mako).Retrosendar bruske}
\l{\cel{3}suflado de aero, produktante bruiso. - EFIS.}
\l{}
\l{\sou{rukular}. (\sou{netrans.}) (\sou{aludante la kolombi, la turturi}).Emisar}
\l{\cel{3}la murmuro karezanta qua karakterizas lia speco. - F.}
\l{}
\l{\sou{rular}. I. (\sou{trans.)}\cel{2}Movar per igar rotacar sur su ipsa. -}
\l{\cel{3}II. (\sou{netrans.}) Movar su per rotaco sur su ipsa. - III.}
\l{\cel{3}(\sou{aludante navo)}. Ocilar de dextre a sinistre, ed inverse).}
\l{\cel{3}- DEF.}
\l{}
\l{\sou{rulado}. (\sou{muziko}).Vokalizo qua konsistas ye intersucedo rapida}
\l{\cel{3}e lejera de noti sur la sama silabo.- EF.}
\l{}
\l{\sou{ruleto}. Hazardo-ludo en qua mikra bulo de ivoro, per rular sur}
\l{\cel{3}cirklo dividita aden 76 faketi numeroza (reda o nigra), de-}
\l{\cel{3}terminas\cel{2}la gano o la perdo, segun ke ta bulo haltas sur la}
\l{\cel{3}faketo. - DEFIRS.}
\l{}
\l{\sou{rumo.}\cel{2}Brandio facita ek la reziduo de melaso di kano-sukro.-}
\l{\cel{3}EFIR.}
\l{}
\l{\sou{rumbo}. (\sou{navig.}) Angulo-spaco qua inter-separas la 32 vent-arei}
\l{\cel{3}di la busolo. - DEFIS.}
\l{}
\l{\sou{ruminar}. (\sou{trans. e netrans}.) (\sou{aludante ula mamiferi herbivora}).}
\l{\cel{3}Igar la manjaji mi-mastikita retrovenar de la stomako aden}
\l{\cel{3}la boko, ube li mastikesas itere ante esar englutata defi-}
\l{\cel{3}nitive. - EFIS.}
\l{}
\l{\sou{rumoro.}\cel{2}Informo-vorti qui propagesas, difuzesas en publiko}
\l{\cel{3}(+publico). - EFIS.}
\l{}
\l{\sou{rumpo}. Extremajo dopa di la korpo di la ucelo, qua suportas}
\l{\cel{3}la kaudo. - E.}
\l{}
\l{\sou{rumsteko}. La parto maxim supra di la dopajo di bovo, an-sur la}
\l{\cel{3}kaudo. - DEF.}
\l{}
\l{\sou{runo.}\cel{1}Singla de la literi ek la alfabeti maxim antiqua di la}
\l{\cel{3}populi Germana e Skandinaviana. - DEFIRS.}
}
